\section{Related Work}
\label{sec:related}

本节从传统 QoS 路由、深度强化学习驱动的路由、图神经网络的网络建模能力、以及图网络条件化机制四个维度，综述与本问题密切相关的近期研究，并在末段明确本文定位。

\subsection{传统 QoS 路由与流量工程}

OSPF、ECMP、CSPF 等经典协议基于静态链路权重做逐跳转发，在混合业务场景下缺乏差异化服务能力。Wang 与 Crowcroft 奠定了 QoS 路由的约束优化框架\cite{Wang1996_QoSRouting}；Fortz 与 Thorup 揭示了基于权重优化的流量工程在动态需求下的局限性\cite{Fortz2000_TrafficEngineering}。Brundiers 等人研究了 Segment Routing 中环路对流量工程的增益，展示了灵活路径控制在动态负载均衡中的潜力\cite{Brundiers2021_Loops_SR_TE}，并进一步提出了中点优化方法以实现更细粒度的流量工程控制\cite{Brundiers2022_Midpoint_SR_INFOCOM}。这些方法的共同缺陷在于短视性:每步决策独立优化，无法在网络全局与序列长远视角下为异构流预留资源。当老鼠流与大象流共享链路时，静态策略往往在保障低时延与追求高吞吐之间顾此失彼。

\subsection{深度强化学习驱动的网络路由}

为突破固定规则的局限，研究者将 DRL 引入路由决策。Liu 等人针对多类型业务需求提出了 DRL-OR，在在线路由中实现了对异构服务的差异化处理\cite{Liu2021_DRLOr_INFOCOM}。Casas-Velasco 等人提出了 DRSIR，以深度 Q 学习在 SDN 中实现智能路由决策\cite{CasasVelasco2021_DRSIR_TNSM}。Chen 等人提出了基于多智能体元强化学习的自适应多路径路由优化框架\cite{Chen2022_MultiagentMeta_TNNLS}。Suzuki 等人提出了多智能体深度强化学习方法，在多云边缘网络中联合优化计算卸载与路由决策，验证了多智能体协作在动态网络环境中的自适应能力\cite{Suzuki2023_MADRL_Routing_TNSM}。Sanchez 等人提出了 DQS，以 QoS 驱动的深度强化学习实现 SDN 中的路由优化\cite{Sanchez2024_DQS_JPDC}。然而，上述方法存在两个共性局限：其一，策略网络大多以全连接层或简单嵌入编码网络状态，缺乏对拓扑结构的显式建模能力，泛化性受限；其二，它们通常为所有流优化单一的全局目标（如最小化最大链路利用率或平均时延），未将逐流 QoS 意图作为条件信号嵌入策略网络，无法在混合业务下实现差异化路由。

\subsection{图神经网络的网络建模与路由优化}

图神经网络（GNN）为编码网络拓扑提供了自然工具。Rusek 等人利用消息传递神经网络建模链路时延与丢包，开创了 RouteNet 系列工作\cite{Rusek2020_JSAC_RouteNet}。Ferriol-Galm\'{e}s 等人进一步提出 RouteNet-Fermi，通过增强的 GNN 架构实现大规模网络性能预测与优化\cite{FerriolGalmes_2023_RouteNetFermi}。Huang 等人提出了xNet，通过学习时序状态转移函数实现细粒度流级 QoS 预测\cite{Huang2024_ToN_xNetModeling}。Munikoti 等人系统综述了深度强化学习与图神经网络结合的挑战与机遇，揭示了现有方法在复杂网络优化中的瓶颈\cite{Munikoti2024_DRL_GNN_TNNLS}。Almasan 等人探索了 GNN 与 DRL 在路由优化中的联合应用，验证了图结构感知对策略质量的提升\cite{Almasan2022_RLmeetsGNN}。Bern\'{a}rdez 等人提出了 MAGNNETO，以基于 GNN 的多智能体系统在流量工程优化中实现了分布式决策与快速响应\cite{Bernardez2023_MAGNNETO_TCCN}。Bhuvanasi 等人利用 GNN 与多智能体 DRL 实现了对拓扑变化的弹性路由\cite{Bhuvanasi2023_ResilientRouting_TNSM}。Yang 等人将图卷积网络与 DRL 结合，在时间敏感网络中联合优化路由与调度\cite{Yang2022_GCN_RL_IoTJ}。Dong 等人提出了基于 GCN 的多任务深度强化学习方法，在网络切片与路由的联合优化中验证了图感知策略的有效性\cite{Dong2022_GCN_Slicing_Routing_TCCN}。然而，现有 GNN 方法存在结构性的表示同质化问题：它们要么聚焦于聚合流量下的离线性能建模（如 RouteNet），学习``一张图、一种全局嵌入''；要么在在线决策中为所有流生成统一的拓扑表示（如 GNN+DRL 方法）。这意味着，当老鼠流与大象流先后到达时，系统看到的是同一张图的同一个嵌入，无法根据逐流的 QoS 特征对同一拓扑进行差异化感知，更无法据此做出差异化的路由决策。

\subsection{图神经网络的条件化机制}

为使网络表示能够根据外部条件动态调整，研究者提出了多种条件化机制。Perez 等人提出的 FiLM（Feature-wise Linear Modulation）通过仿射变换将辅助信息注入神经网络各层 \cite{Perez2018_FiLM}；Brockschmidt 等人将其拓展至图领域，提出 GNN-FiLM，以目标节点表示调制消息传递过程 \cite{Brockschmidt2020_GNNFiLM}。在异构图学习领域，Hu 等人提出的 Heterogeneous Graph Transformer 通过节点类型条件化注意力来区分不同语义关系 \cite{Hu2020_HGT}；Brody 等人提出的 GATv2 以动态注意力机制实现了边级条件化权重计算 \cite{Brody2022_GATv2}。然而，上述条件化机制主要服务于节点分类、分子生成或静态图表示学习任务，其调制信号来源于图内结构（目标节点、边类型、节点类别等），而非逐流的外部 QoS 意图。在在线路由场景中，真正需要条件化的并非图内节点属性，而是到达流的 QoS 特征：同一张拓扑在面对老鼠流与大象流时，应当产生截然不同的嵌入与决策。现有文献尚未将 FiLM 条件化机制引入逐流在线路由场景，使得混合业务下的差异化拓扑感知与策略学习成为开放问题。


综上，现有工作存在明确缺口：（1）传统方法短视且无法处理混合 QoS；（2）DRL 方法缺乏拓扑显式建模；（3）GNN 方法假设同质流量；（4）条件化 GNN 尚未引入逐流在线路由。本文填补上述差距，提出流感知 FiLM-GNN 架构，以逐流 QoS 意图作为外部条件信号，通过 FiLM 重塑 GNN 的每层消息传递，使同一张物理拓扑在异构业务流下被编码为差异化的策略输入，进而以 PPO 训练非短视的在线选路策略。